%%===================================================================%%
%% Results — revision of 2026-09-21                                  %%
%% Supersedes results_body.tex (2026-09-19).                         %%
%% Every number here is traceable to a file under results/ produced  %%
%% by rerun_all.sh on the 796 / 919 analysis sets.                   %%
%%===================================================================%%

\section{Results}

\subsection{Temporal degradation in both datasets}

The analysis sets comprise 796 DILIrank and 919 DICTrank compounds with an
approval year, a standardised structure and a binary label. Splitting at 2010
leaves 640 training and 156 held-out compounds for DILIrank, and 652 and 267 for
DICTrank.

On the DILIrank held-out set the model reached an AUROC of 0.648 (95\% bootstrap
CI 0.555--0.736). One thousand random partitions of the same held-out size
averaged 0.735 (SD 0.036, 2.5--97.5 percentile 0.663--0.806), and six of them fell
at or below the observed value. Using the bias-corrected estimator of Phipson and
Smyth, which cannot return zero, the one-sided permutation $p$ is 0.0070.
DICTrank behaved the same way: 0.620 (0.553--0.688) against a random reference of
0.697 (SD 0.029, 0.636--0.752), four partitions at or below, $p = 0.0050$. Both
survive Benjamini--Hochberg correction within the confirmatory family at
$q = 0.0070$.

The held-out sets are also less positive than the training sets, so part of the
AUROC gap could be an artefact of class balance. Resampling the DILIrank held-out
set to the training positive rate of 65.5\%, averaged over 20 draws, moved AUROC
from 0.648 to 0.632 on the resulting 73 compounds --- that is, downwards. Measured
against the 20-partition random control used in the same script, which gives
0.742, the gap widens from 0.094 to 0.110 rather than closing. The direction
matters: class
balance is not masking a smaller gap, and the shift accounts for none of the
degradation in discrimination. Its effect on calibration is a separate question
and is treated below.

Gaps of 0.09 and 0.08 are modest. Reported alone they would say only that
prospective validation is harder than retrospective validation, which is already
well known. The remainder of this section concerns what those numbers conceal.

\subsection{The degradation is a calibration loss}

Table~\ref{tab:brier} gives the Murphy decomposition. Each random partition is
decomposed at the same held-out size as the temporal split and the terms are then
averaged, so the two columns rest on equal sample sizes.

For DILIrank the Brier score rose from 0.213 to 0.427. Resolution, which tracks
discrimination, fell by 0.034. Reliability, which tracks calibration, rose by
0.208. Calibration therefore accounts for \textbf{86.1\%} (95\% CI 84.6--97.3) of
the degradation against 13.9\% for discrimination, a ratio of 6.2 to 1. DICTrank
gives \textbf{82.5\%} (75.5--124.0) against 17.5\%, a ratio of 4.7 to 1: a Brier
score of 0.205 rising to 0.306, resolution down 0.012, reliability up 0.054.

Pooling the random partitions before binning, as an earlier version of this
analysis did, inflates both shares --- to 90.4\% for DILIrank and 91.8\% for
DICTrank. Pooling gives the random reference twenty times the sample of the
temporal column, which shrinks its reliability term and flatters the comparison by
4.3 and 9.3 percentage points. The per-partition figures are reported as primary
throughout.

The ordinary calibration metrics agree. Against a random split of the same size,
expected calibration error rose from 0.121 to 0.438 in DILIrank and from 0.126 to
0.270 in DICTrank, and maximum calibration error from 0.295 to 0.622 and from
0.471 to 0.593. The calibration slope fell from 0.534 to 0.270 and from 0.433 to
0.222, and the intercept moved from $+0.053$ to $-1.313$ and from $+0.326$ to
$-0.085$. Mean predicted probability on the DILIrank held-out set was 0.732 where
the observed positive rate was 0.308. For DICTrank the figures are 0.845 and
0.614.

These are over-confident predictions, and the over-confidence is worst in the bins
a user would act on. Seventy-one of the 156 DILIrank compounds were assigned a
probability above 0.9, with a mean prediction of 0.962; 40.8\% of them were
positive. The corresponding DICTrank bin held 156 of 267 compounds at a mean
prediction of 0.966 against an observed 67.9\%. Figure~\ref{fig:calib} shows the
reliability curves in full.

Because the failure is a shift rather than a scrambling, moving the decision
threshold recovers part of what is lost at 0.5. The threshold maximising Youden's
J was 0.920 for DILIrank and 0.940 for DICTrank. At those thresholds balanced
accuracy rose from 0.557 to 0.635 and from 0.519 to 0.601, and the Matthews
correlation coefficient from 0.126 to 0.254 and from 0.062 to 0.198. A threshold
near 0.93 is not a usable default; the point is that one has to move that far to
find it.

AUROC is invariant to any monotone transformation of the predicted probabilities,
so none of this appears in it.

\subsection{Scaffold splitting does not reproduce the calibration collapse}

Compounds approved after 2010 are structurally newer as well as later, which
raises the possibility that any sufficiently hard split would produce the same
picture. Running stratified and Bemis--Murcko scaffold cross-validation on the
same compounds, model and features tests that directly
(Table~\ref{tab:split}, Figure~\ref{fig:split}a,b).

Discrimination gave an inconsistent answer. Scaffold cross-validation raised
DILIrank AUROC slightly, from 0.728 $\pm$ 0.022 to 0.731 $\pm$ 0.018, and lowered
DICTrank from 0.699 $\pm$ 0.023 to 0.678 $\pm$ 0.036. Both changes are smaller
than the fold-to-fold standard deviation, and they point in opposite directions in
the two datasets. Read through AUROC alone, scaffold splitting is neither reliably
harder nor reliably easier than a random split here.

Calibration gave a consistent one. Scaffold ECE was 0.130 $\pm$ 0.029 against
0.136 $\pm$ 0.012 for DILIrank and 0.136 $\pm$ 0.030 against 0.129 $\pm$ 0.022 for
DICTrank --- changes of $-0.006$ and $+0.007$, within a standard deviation in both
cases. The temporal split on the same compounds gave 0.438 and 0.270, which is
3.4 times the scaffold value in DILIrank and 2.0 times it in DICTrank. Structural
novelty in the held-out set, imposed deliberately, leaves calibration where a
random split leaves it; temporal ordering does not.

One reason the two cross-validation designs resemble each other is visible in the
scaffold counts. The 796 DILIrank compounds contain 542 unique Bemis--Murcko
scaffolds, 455 of which appear exactly once; the 919 DICTrank compounds contain
638, with 545 singletons. When 84\% and 85\% of scaffolds are singletons,
grouping by scaffold removes little that a random partition would not already
separate, and the design converges on stratified cross-validation. This is a
property of drug-like reference sets rather than of these two datasets in
particular, and it is worth stating plainly before scaffold splitting is offered
as a proxy for prospective validation.

This comparison was not among the pre-specified hypotheses, and it is reported as
a description rather than a test. No $p$-value is attached to it.

\subsection{Prevalence-matched controls}

The positive rate falls across the split boundary, from 65.5\% to 30.8\% in
DILIrank and from 77.0\% to 61.4\% in DICTrank, and a model calibrated to one base
rate will be miscalibrated on another. Three nested controls, each averaged over
200 partitions, separate that effect from the rest (Figure~\ref{fig:split}c).
Because they use a different number of repeats from the calibration analysis
above, their random-split values differ from it in the second decimal place.

DILIrank temporal ECE was 0.438 (95\% CI 0.374--0.527). An unconstrained random
split (R1) gave 0.131 $\pm$ 0.029, a ratio of 3.35. Matching the held-out positive
rate (R2) raised the random-split value to 0.250 $\pm$ 0.032, ratio 1.76. Matching
both rates (R3) raised it to 0.321 $\pm$ 0.032, ratio 1.37. DICTrank temporal ECE
was 0.270 (0.228--0.336), against 0.130, 0.187 and 0.223 for ratios of 2.08, 1.45
and 1.21.

No random partition out of 200 exceeded the temporal ECE in five of the six
comparisons, giving $p = 0.0050$, the lower bound available at 200 repeats under
the Phipson--Smyth estimator. The exception is DICTrank R3, where six of 200 did,
giving $p = 0.0348$. All six survive correction within their family, the largest
$q$ being 0.0348.

Matching both base rates closes most but not all of the excess: of the 0.308 of
ECE separating the temporal split from an unconstrained random one in DILIrank,
0.117 survives R3, a closure of 62\%; of the 0.140 in DICTrank, 0.048 survives, a
closure of 66\%. Prevalence shift is a large part of the calibration collapse and
not the whole of it, and the two datasets agree closely on how large a part.

Discrimination behaves differently under the same controls. R3 AUROC was 0.726 for
DILIrank and 0.693 for DICTrank against temporal values of 0.648 and 0.620, so
matching the base rates removes little of the AUROC gap while removing most of the
ECE gap. In DILIrank the fully matched comparison is a 10.8\% relative loss in
AUROC against a 36.5\% relative rise in ECE.

\subsection{Recalibration under a prospective constraint}

Holding out the most recent quarter of the training compounds, ordered by approval
year, as an internal calibration set keeps strategies (B) and (C) inside the
information a prospective user would have; strategy (D) does not, for the reason
given below. Table~\ref{tab:recal} gives the result, and it is a mixed one.

DILIrank began at ECE 0.447. The oracle ceiling, Platt scaling fitted on the
held-out labels, reached 0.048. Internal Platt scaling closed 8\% of that gap and
isotonic regression 10\%. Only the prior-probability correction moved the needle,
reaching 0.296 and closing 38\%. DICTrank reversed the ordering: internal Platt
scaling reached 0.168 and closed 42\%, isotonic regression 0.174 and 40\%, and the
prior correction 0.244 and 12\% against an oracle of 0.023.

The prior-probability correction is not fully prospective and is reported with
that caveat. It rescales the intercept using the held-out positive rate, which a
user would have to estimate rather than know. It is included because it isolates
how much of the miscalibration is a pure base-rate effect; strategies (B) and (C)
are the ones that use no held-out information at all.

Neither of the two fully prospective strategies recovered even half of the
recoverable error, and the two datasets differ by a factor of four to five in how
much they recover at all: 8--10\% for DILIrank against 40--42\% for DICTrank. The
base-rate correction runs the other way, recovering 38\% in DILIrank and 12\% in
DICTrank. The positive rate moves by 35 percentage points across the DILIrank
split boundary and by 16 across the DICTrank one, which is the ordering the
base-rate correction follows. The prevalence-matched controls do not follow it ---
they close a similar fraction of the ECE excess in both datasets, 62\% and 66\% ---
so the two analyses are answering different questions and should not be read as
corroborating each other.

Recalibration left discrimination where it was, as monotone transformations must:
AUROC was unchanged to four decimal places under the prior correction. The 0.08 to
0.09 of lost AUROC is a separate problem, and nothing here addresses it.

One bookkeeping point. The uncorrected condition in this analysis gives AUROC
0.660 and ECE 0.447 for DILIrank, against 0.648 and 0.438 in the primary split.
The recalibration pipeline sorts compounds by approval year before fitting, which
changes which rows XGBoost draws under row subsampling. Both values fall inside
the ten-seed range reported below, and every recovery fraction in
Table~\ref{tab:recal} is computed against the uncorrected condition of this same
pipeline.

\subsection{Conformal coverage}

Split conformal prediction offers a distribution-free coverage guarantee that
rests on exchangeability, which a temporal split violates by construction. The
guarantee fails; the interesting part is how (Table~\ref{tab:conformal}).

Marginal coverage fell in both datasets and further in DILIrank. At a nominal
90\%, the DILIrank random reference delivered 90.7\% and the temporal split
73.7\%; at 80\%, 81.0\% against 48.1\%. DICTrank lost less: 89.9\% against 82.4\%
at 90\% nominal, and 80.7\% against 69.7\% at 80\%. That ordering matches the
ordering of the ECE ratios reported above, and is obtained without any
distributional assumption.

Splitting coverage by class shows where the loss falls, and here DICTrank is the
more extreme of the two. Under the DILIrank temporal split the positive class was
covered at 95.8\%, above nominal, while the negative class reached 63.9\% --- a
separation of 32 percentage points, widening to 63 points at nominal 80\% (91.7\%
against 28.7\%). DICTrank separated further still: 99.4\% against 55.3\% at 90\%
nominal, 44 points, and 98.8\% against 23.3\% at 80\%, 75 points. So DICTrank is
the milder case by marginal coverage and the more severe one by class balance,
and reporting only the marginal figure would reverse the impression.

The shortfall falls almost entirely on the negative class in both datasets. In
this setting the negative label is the reassuring one, so a prediction set that
omits it is a model refusing to certify as safe a compound that the label calls
safe. Marginal coverage alone reports a guarantee that is somewhat violated and
hides which half of the guarantee has gone.

Empty prediction sets occurred in one condition only --- 0.3\% of the DICTrank
random reference at nominal 80\% --- and in none of the temporal conditions, so
the shortfall is not an artefact of the quantile running off the end of the
calibration scores. Split conformal guarantees marginal coverage only, so the
class-conditional figures are descriptive; this is why the random reference itself
falls below nominal within the negative class, most visibly at 48.4\% for DICTrank
at nominal 80\%.

\subsection{Label maturity}

Post-marketing labels accumulate, and held-out compounds are younger, so some of
the negative labels in the held-out set may be negative only for now. A gradient
is present in the data over the first 25 years of observation. Among DILIrank
compounds with no more than seven years of observation, 14.0\% carry a positive
label; this rises to 28.6\% at 8--11 years, 52.0\% at 12--15 years and 67.7\% at
16--25 years. DICTrank shows a flatter version of the same trend, from 59.3\% to
78.4\% over the same range.

The gradient does not continue indefinitely. In the oldest stratum, compounds with
more than 25 years of observation --- which is 516 of 796 DILIrank compounds and
472 of 919 DICTrank compounds, the majority of each dataset --- the positive rate
falls back to 64.9\% and 76.1\%. Label accumulation therefore appears to saturate,
and the trend should not be extrapolated as though every negative label will
eventually turn positive. The DICTrank strata sum to 916 rather than 919: three
compounds --- nilotinib, amantadine and selumetinib --- carry an approval year
later than the 2023 label release, which the audit flags as a probable
approval-year matching error rather than a genuine post-release approval. They
have no defined observation window and are therefore absent from this analysis
only. Two of the three, nilotinib and amantadine, are among the corrections tested
in the approval-year sensitivity analysis below; selumetinib is not, because its
earliest matching record is also after 2010 and correcting it would not move it
across the boundary.

Restricting the held-out set to mature labels narrows the DILIrank gap without
closing it. Against a baseline gap of 0.094, requiring at least eight years of
observation gave 0.084, at least ten years 0.077, and at least twelve years 0.072
--- a reduction of 24\% on 50 held-out compounds. Down-weighting negative labels
in proportion to observation time changed the gap by less than 0.003 at every
value of $\tau$. Applying the restriction to training and held-out sets together
widened the gap, to 0.106.

DICTrank did not move in the expected direction at all. Against a baseline of
0.088, the eight-year restriction widened the gap to 0.117 and the ten-year
restriction to 0.093. The largest reduction available under any control was
0.0020, from down-weighting at $\tau = 5$, and the joint restriction widened the
gap to 0.111.

Label maturity therefore accounts for at most about a quarter of the
discrimination gap in DILIrank, on the smallest and most favourable subset
examined, and for essentially none of it in DICTrank. In both datasets the
comparison that holds training and held-out maturity to the same standard widens
rather than narrows the gap. These four controls were evaluated on AUROC only;
calibration was not recomputed under them, so they bear on the discrimination gap
and not on the calibration collapse.

\subsection{Robustness}

Across all 16 combinations of four learners and four representations
(Figure~\ref{fig:robust}a), ECE was worse under the temporal split in
\textbf{16 of 16} cells for both datasets. The median ratio was 2.64 for DILIrank
(range 1.64--5.28) and 1.91 for DICTrank (range 1.43--3.65), and the median
calibration slope 0.184 and 0.172. Discrimination was less uniform. $\Delta$AUROC
was positive in 15 of 16 DILIrank cells, median $+0.096$, and in 16 of 16 DICTrank
cells, median $+0.088$. The one exception, RDKit descriptors with logistic
regression, is also among the weakest cells in the grid under random splitting, so
there was little discrimination there to lose.

Moving the cut-off to 2005 or 2015 preserves the calibration pattern and not the
discrimination pattern (Figure~\ref{fig:robust}b). ECE ratios stayed between 2.71
and 3.67 for DILIrank and between 1.64 and 2.19 for DICTrank, never falling below
1.6 in either dataset. $\Delta$AUROC was the less stable quantity, ranging from
0.041 to 0.097 for DILIrank --- a factor of 2.4 --- and from 0.050 to 0.082 for
DICTrank. Both datasets peak at 2010, but the 2005 and 2015 values swap rank
between them, so the size of the AUROC gap is not a stable function of how far
into the past the boundary is placed. The reliability
share of the Brier degradation was 88.2\%, 86.1\% and 81.4\% for DILIrank at 2005,
2010 and 2015, and 72.9\%, 82.5\% and 90.1\% for DICTrank. The share exceeds 70\%
at every cut-off in both datasets, but it is not constant --- it spans about 7
percentage points in DILIrank and 17 in DICTrank --- and the two datasets move in
opposite directions as the cut-off is made later.

Ten seeds moved AUROC by a standard deviation of 0.0094 for DILIrank and 0.0103
for DICTrank, and ECE by 0.0071 and 0.0074; the seed-42 values reported throughout
sit inside both ranges. The DILIrank AUROC of 0.648 is close to the ten-seed mean
of 0.647; the DICTrank value of 0.620 sits in the upper part of its range
(0.591--0.626, mean 0.611), which is the conservative direction for the claim
being made. Recomputing ECE and the decomposition over 5, 10, 15 and 20 bins under
equal-width and equal-frequency binning moved ECE between 0.424 and 0.454 for
DILIrank and between 0.231 and 0.282 for DICTrank. Reliability exceeded resolution
under every scheme, by a margin of at least 0.19 for DILIrank and 0.04 for
DICTrank.

\subsection{Sensitivity to data-curation decisions}

Three curation decisions could each be argued to drive the result, and each was
tested directly.

\textbf{Approval-year assignment.} An audit against Drugs@FDA identified seven
compounds in the analysis sets assigned to the held-out side of the boundary whose
earliest matching approval record predates 2010: one in DILIrank and six in
DICTrank. Applying those corrections (S1), and additionally excluding phenobarbital
sodium, in clinical use since 1912 but first recorded as an NDA in
2022 (S2), moved the DILIrank ECE ratio from 3.37 to between 3.37 and 3.68 and the
reliability share from 86.1\% to between 84.7\% and 86.2\%. The DICTrank ratio
moved from 1.96 to 2.05 and the share from 82.5\% to 80.3\%. The largest movement
in either dataset is 0.31 in the ratio and 2.2 percentage points in the share.

\textbf{Diagnostic and contrast agents.} Compounds in ATC class V carry safety
labels of a different character from therapeutic agents. Removing them (44
compounds from each dataset, leaving 752 and 875) left the DILIrank ECE ratio at
3.44 against 3.37 and the DICTrank ratio at 1.85 against 1.96, and raised the
reliability share in both datasets, to 91.2\% and 86.7\%. The calibration collapse
is not an artefact of this class.

\textbf{Structure-based deduplication.} Collapsing salts and stereoisomers to a
common parent structure merges entries that occasionally carry opposite labels ---
for instance vancomycin and vancomycin hydrochloride in DICTrank. Rather than
retaining an arbitrary member of such a group, every conflicted group was removed
in full: one group in DILIrank and thirteen in DICTrank, which removes one entry
from the DILIrank analysis set and thirteen from the DICTrank one.
The rule was fixed before its effect was examined, and its effect was not
one-directional: re-running the decomposition on the retained pre-filter sets
(932 and 797 compounds) gives a reliability share of 85.6\% for DILIrank against
86.1\% after the rule, and 84.5\% for DICTrank against 82.5\% --- a rise in one
dataset and a fall in the other. A filter chosen to favour the reported result
would not behave this way.
The pre-filter analysis sets, the re-run output and the full list of removed
groups are given in Additional file~1.

\subsection{Exploratory subgroup analyses}

Held-out compounds sit far from the training set in chemical space. The median
nearest-neighbour Tanimoto similarity to training compounds was 0.271 for
DILIrank, and 115 of 156 held-out compounds (73.7\%) fell below 0.4. Performance
rose with similarity: AUROC 0.543 below 0.3 ($n = 92$), 0.711 between 0.3 and 0.4
($n = 23$), and 0.857 between 0.4 and 0.6 ($n = 29$). The formal tests do not
support the trend. Splitting the held-out set at the median gave a difference of
$+0.068$ with a two-sample bootstrap $p$ of 0.527, and the correlation between
similarity and the predicted probability of the true class was weak
(Spearman $\rho = +0.116$, $p = 0.148$).

The largest subgroup effect was therapeutic rather than structural. Antineoplastic
and immunomodulating agents (ATC class L) accounted for 32 held-out DILIrank
compounds, 84\% of them positive, with an AUROC of 0.289 against 0.649 for the
remaining 118, a difference of $-0.360$ at a bootstrap $p$ of 0.021
(Figure~\ref{fig:atc}). The same comparison in DICTrank gave $-0.048$,
$p = 0.552$, on a subgroup three times larger. Of the eight physicochemical
strata, two reached nominal significance --- DICTrank molecular weight above 500
($p = 0.020$) and DICTrank six or more nitrogen atoms ($p = 0.050$) --- and the
corresponding DILIrank nitrogen stratum points the opposite way.

After Benjamini--Hochberg correction within the exploratory family,
\textbf{none of the twelve tests remains significant}; the smallest $q$ is 0.126,
carried jointly by the DILIrank ATC class L result and the DICTrank molecular
weight stratum. Below-chance AUROC in a 32-compound subgroup is reported for
completeness and supports no claim made in this paper.
